This work presents a unified benchmark evaluation of RAG and GraphRAG across both question answering and query-based summarization, clarifying when explicit graph structures help, and when they do not, under controlled settings. Our analyses reveal strong task-dependent behaviors: RAG is consistently effective for single-hop, detail-oriented queries that require precise evidence, whereas GraphRAG is more advantageous for multi-hop, reasoning-intensive QA and tends to produce more corpus-level, diverse summaries.
Motivated by these findings, we study two hybrid strategies, \textbf{Selection} and \textbf{Integration}, that combine the strengths of both paradigms and improve QA performance. Beyond effectiveness, we highlight practical challenges that limit current GraphRAG systems, including incomplete or noisy graph construction, additional computation and storage overhead, and evaluation artifacts such as position effects in LLM-as-a-Judge for summarization. Together, these observations point toward the next generation of RAG systems: approaches that can construct and refine graphs reliably, adapt retrieval and aggregation to query needs, and deliver stronger reasoning benefits under realistic efficiency constraints.
